%! Project: Design and Conduct a Survey — Unit 1, Lesson 1.8 — Lesson Plan
% PILOT SHAPE (2026-09-06): modeled on AP Statistics lesson 1.4 minus its AP Practice page.
% GRADUAL-RELEASE plan (warm-up / guided notes and practice / spoken debrief / close and START
% THE HOMEWORK). The guided notes carry the release ROW BY ROW in one Main Ideas / Notes table
% (2026-09-12 shape): I do / we do in rows 1-4 (problems 1-15) -> we do the Guided Practice row
% (problems 16-19), which is the class's OWN survey, run live. THERE IS NO INDEPENDENT PRACTICE
% SET IN THE NOTES — the homework is the individual practice, and its item 6 is the project
% deliverable. Teacher-facing. Breakable boxes are `skillbox` with a \boxguard ahead; the two
% tabularx boxes are `fixedskillbox` (never splits). No group activity, no tier boxes, no exit
% ticket — the old exit ticket's Bayside homeroom is now the debrief's cold check.
\documentclass[10pt]{article}
\usepackage{saar-article}
\usepackage{saar-boxes}
\usepackage{graphicx}   % -article does NOT load graphicx
\graphicspath{{images/}}
\newcommand{\TallMath}[1]{$\displaystyle #1\rule[-1.4em]{0pt}{3.2em}$}
\newcommand{\UnitNumberName}{Unit 1: Asking Questions with Data: Study Design \quad}
\newcommand{\LessonNumberName}{Lesson 1.8: Project: Design and Conduct a Survey}

\begin{document}
\begin{center}
    {\Large\bfseries \CourseName \\
    {\small\bfseries \UnitNumberName \LessonNumberName}}
\end{center}

\begin{tcolorbox}[colback=frost, colframe=cerulean, boxrule=0.9pt, arc=2mm,
    left=2mm, right=2mm, top=1.5mm, bottom=1.5mm]
    \textbf{Primary Objective:} Students will run the four-stage statistical cycle forwards on a
    survey of their own --- write a statistical question, judge a sampling plan, build a
    \emph{survey instrument} whose items ask one thing in neutral words with choices that cover
    everyone, tally real answers from real classmates into relative frequencies, read them off a
    pre-drawn bar graph, and write a four-sentence report --- and will judge that report's last
    sentence by naming whom the design actually allows them to describe.
    \\[2pt]
    \textbf{Standards:} AFDA.DA.2d, AFDA.DA.2h, AFDA.DA.2i, AFDA.DA.2j; PS.DC.2d
    \quad$\cdot$\quad AFDA.DA.2a--c and 2e--g are exercised in the planning rather than taught
    fresh \quad$\cdot$\quad closes Unit 1 (PS.DC.1--2 run forwards)
    \\[2pt]
    \textbf{Lesson model:} \emph{Gradual release --- I do, we do, you do --- all of it inside the
    guided notes.} There is no group activity. The notes name the vocabulary as they teach it and
    deliver the instruction \textbf{row by row in one Main Ideas / Notes table} (four rows,
    problems 1--15: you read the definition, mark up the display, and work the first problem of
    each grid; the class works the rest); \textbf{the Guided Practice row (problems 16--19) is the
    class's own survey}, built and tallied live with students holding the pen; \textbf{the homework
    is the individual practice}, and students start it in class while you circulate. The whole-class
    debrief is \emph{spoken}.
    \\[2pt]
    \textbf{Note --- this is the unit's project lesson, and it is one lesson.} The project is not a
    separate week: the instrument is built in the Guided Practice row, the data are collected in the
    room in about four minutes, and the deliverable is \textbf{homework item~6}, the four-sentence
    report on the class's own tally. Decide \emph{before} class which single question the room will
    vote on, so the four minutes are spent asking, not negotiating.
\end{tcolorbox}

\renewcommand{\arraystretch}{1.4}
\boxguard
\begin{skillbox}[Learning Targets \& Key Understandings]{goldbox}
\begin{tabularx}{\linewidth}{@{}X|X@{}}
\textbf{Learning targets (student language)} & \textbf{Key understandings (the ``why'')} \\[2pt]
\hline
\vspace{-6pt}
\begin{itemize}[leftmargin=1.1em, nosep, after=\vspace{2pt}]
  \item I can plan a survey through the four stages of the \textbf{statistical cycle} and say
        what each stage hands in. \hfill \textit{PS.DC.2d}
  \item I can write a \textbf{survey instrument} whose \textbf{closed-response items} ask one
        thing, in neutral words, with choices that cover everyone without overlapping.
        \hfill \textit{AFDA.DA.2d}
  \item I can turn a tally into \textbf{relative frequencies} and read them off a \textbf{bar
        graph}. \hfill \textit{AFDA.DA.2h--i}
  \item I can report a conclusion about the people my design actually covers, and explain why a
        \textbf{convenience sample} that happens to agree with a chance sample has proved
        nothing. \hfill \textit{AFDA.DA.2j}
\end{itemize}
&
\vspace{-6pt}
\begin{itemize}[leftmargin=1.1em, nosep, after=\vspace{2pt}]
  \item A survey is an \emph{observational study}, so everything Unit~1 built about somebody
        else's study applies to the students' own.
  \item Each stage hands something to the next. \emph{A bad question cannot be repaired by good
        arithmetic later}, and bias written into the instrument is in the data permanently.
  \item \textbf{Target misconception:} that near-agreement is evidence. Last spring's class of
        $25$ got $36\%$ against the council's $35\%$, and half the room will read that as
        proof that a class can speak for the school. \emph{Chance never chose the class}, so the
        match is luck --- the same room could have missed by twenty points and nothing in its
        method would have told anyone.
  \item \textbf{Second misconception:} that correct arithmetic launders a bad sample.
        $0.36 \times 900 = 324$ is right and means nothing; $324$ describes nobody.
  \item Naming the limit is \emph{part of the answer}, not an apology --- it is what
        AFDA.DA.2j scores.
\end{itemize}
\end{tabularx}
\end{skillbox}

\boxguard[24]
\begin{skillbox}[{Vocabulary, Concepts, \& Theorems --- taught directly in the guided notes}]{greenbox}
\small
\renewcommand{\arraystretch}{1.35}
\begin{tabularx}{\linewidth}{@{} >{\bfseries}p{3.1cm} X @{}}
Survey instrument & The exact wording read to every person in the sample, in order, plus its
       choices --- the tool that turns a statistical question into data. \\
Closed-response item & A question offering a fixed set of choices, so every answer lands in one
       category and the answers can be counted. \\
Relative frequency & A category's count divided by the total number of responses, as a decimal or
       a percent. The categories' relative frequencies add to $100\%$. \\
Bar graph & The display for categorical data: one bar per category, bars separated, height equal
       to the frequency or the relative frequency. \\
Convenience sample & A sample chosen because it was available --- a class surveying itself, every
       time. Chance had no turn, so the result describes only the people in it. \\
Four instrument rules & One idea per item $\cdot$ neutral wording $\cdot$ choices cover everyone
       without overlapping $\cdot$ short. \emph{(The 1.4 names for breaking them: double-barrelled,
       leading, overlapping choices, too long.)} \\
\end{tabularx}
\end{skillbox}

% fixedskillbox never splits, so the phase table stays intact on one page.
% THE PHASE TABLE IS A CONTRACT: 5 / 35 / 10 / 10 — the I do is ~20 of the 35 and the Guided
% Practice ~15; the debrief is spoken; the homework start is the second formative read.
\begin{fixedskillbox}[Lesson at a Glance --- \MeetingLength]{frost}
\renewcommand{\arraystretch}{1.25}
\begin{tabularx}{\linewidth}{@{}l c X X@{}}
\textbf{Phase} & \textbf{Min} & \textbf{Students} & \textbf{Teacher} \\
\hline
Warm-Up & 5 & The council's finished survey: $21/60$, $0.35 \times 900$, and the two groups
  hiding in the paragraph --- silent & Debrief item~3 aloud; write ``$60$ asked, $900$ reported
  on'' on the board and leave it up all period \\
Guided Notes \& Practice & 35 & Fill the vocabulary box as each term is named; work rows 1--4
  (problems 1--15) with the pen in their hand; then \emph{Our Class Survey} (16--19) --- build
  the instrument, answer it, tally it & \textbf{I do / we do, row by row}: read the definition,
  mark up the display, work problems 1, 5, 8, 12; the class works the rest (20) $\to$ \textbf{we
  do} the class's own survey, 16--19 (15) \\
Debrief & 10 & Pens down, papers up --- answer aloud, nothing to fill in & Put the council's $60$
  and the class's $25$ back up; take the ``we matched, so we count'' reflex, then the
  ``the arithmetic is right'' reflex; cold check on Bayside \\
Close \& Assign & 10 & \textbf{Start the homework in class}, alone and silent & Announce the
  homework and its form (packet or Desmos), circulate for the second formative read, preview
  the unit review and test \\
\end{tabularx}
\end{fixedskillbox}

\begin{fixedskillbox}[Warm-Up --- Activate Prior Knowledge (5 min)]{frost}
\begin{tabularx}{\linewidth}{@{}X X@{}}
\begin{minipage}[t]{\linewidth}\vspace{0pt}
    \textbf{The three items and what each seeds}
    \begin{itemize}[leftmargin=1.1em, itemsep=1pt, topsep=2pt]
      \item \textbf{1.} $21$ of $60$ as a percent. \textbf{Seeds:} the relative frequency
            column of problem~7 --- the $35\%$ is already on the page when row~3 asks for it.
      \item \textbf{2.} $0.35 \times 900 = 315$. \textbf{Seeds:} scaling a sample percent to a
            population --- problem~8 repeats the move legally, problem~9 repeats it illegally.
      \item \textbf{3.} Two groups, no vocabulary, then the $100\%$ check. \textbf{Seeds:}
            population and sample for rows 1 and~4, and \emph{instrument rule~3} --- one box
            per person is \emph{why} the percents add to $100\%$ (problem~11).
    \end{itemize}
\end{minipage}
&
\begin{minipage}[t]{\linewidth}\vspace{0pt}
    \textbf{Running it}
    \begin{itemize}[leftmargin=1.1em, itemsep=1pt, topsep=2pt]
      \item \textbf{Debrief item~3 aloud} and write the two counts on the board --- \emph{$60$
            asked, $900$ reported on}. Leave it up all period; problem~16 hangs the class's own
            \emph{$25$ asked, $900$ reported on} underneath it, and the debrief circles both.
      \item Item~3's last question is the tell. A student who answers \emph{exactly one} has the
            reason the percents total $100\%$ and will fly through rows 2--3.
      \item Items 1--2 are Lesson 1.3's arithmetic --- say so, so nobody thinks they are being
            retested. Do not let them expand; five minutes, silent, timed.
    \end{itemize}
\end{minipage}
\end{tabularx}
\end{fixedskillbox}

\boxguard
\begin{skillbox}[Hook --- before anyone sits down]{frost}
The hook is on \textbf{page 1 of the notes}, under the vocabulary box, and on the board:
\textit{``Last spring a class of $25$ in this room asked itself the council's grant question and
got $36\%$ for club funding. The council's stratified sample of $60$ had said $35\%$. One student
wrote: `We matched the council, so our class can speak for all $900$ students too.'\,''}
\textbf{Ask: is that student right?} Students circle \emph{agree} or \emph{disagree} and write one
line of reason in the hook box \emph{before} anyone speaks; then take the hand vote and write the
split on the board. Most rooms vote \emph{agree} --- the match looks like confirmation.
\textbf{Do not settle it.} Say only that the answer is problem~13, in \emph{The Limit} row, and
that they will vote again there. The vote is what makes the crux land: students have to have
committed to \emph{matching proves something} before the notes take it away.
\end{skillbox}

\boxguard
\begin{skillbox}[Guided Notes \& Practice --- I do, then we do (35 min)]{frost}
\textbf{This is the whole lesson, and it is where the project lives.} There is no group activity
and no independent practice set on this page: \textbf{the homework is the individual practice},
started in class, and its item~6 is the project deliverable. The notes are \textbf{one Main Ideas
/ Notes table}: four instruction rows (problems 1--15) on the Student Council's \emph{finished}
survey, then the Guided Practice row (16--19), which is the survey \emph{this class runs on
itself}. \textbf{The rhythm in every row is the same}: read the printed sentences aloud, mark up
the display, work the first problem of the grid yourself, then hand the rest to the room. The
vocabulary box is filled \emph{as each term is named}. The only blanks are the two fill tables
(the three broken items in row~2 and the relative frequency column in problem~7) and the
instrument card and tally table of the Guided Practice row.
\begin{multicols}{2}
\textbf{I do / we do, row by row (about 20 min).}

\emph{Row 1, The Cycle (problems 1--4).} Open on the line already on the board: $60$ asked, $900$
reported on. Read the four deliverables, then have students name the stage under each box of the
display. Work problem~1 yourself; the class does 2--4. Problem~3 is Lesson 1.3's stratified
allocation reused unchanged ($270/15 = 18$, $180/15 = 12$); \textbf{problem~4 is the first time
the class is named as a convenience sample} --- take it slowly, it is the seed of the crux.
Four minutes.

\emph{Row 2, The Instrument (problems 5--6) --- the only genuinely new content today.} Read the
four rules, then the card, then the three broken items. Take the fill table from the class; every
one of the three is an item they met in Lesson 1.4 as response bias, so ask for the 1.4 name
\emph{after} they have named the rule. Work problem~5 yourself. Five minutes.

\emph{Row 3, Relative Frequency (problems 7--11).} Problem~7 is warm-up arithmetic; take it fast.
The bar graph is pre-drawn --- nobody sketches anything --- so spend the time on what a bar's
\emph{height} is. \textbf{Problem~9 is the trap}: $9/25 = 36\%$ and $0.36 \times 900 = 324$ are
both correct, and the second number is still not allowed. \textbf{Take a hand count before anyone
writes.} Problem~10 ($42$ of $120$ is still $35\%$) is why a height is a share. Six minutes.

\emph{Row 4, The Limit (problems 12--15) --- the crux.} Put the two panels up, read the
four-sentence shape, work problem~12 live, and then \textbf{re-take the hook vote at problem~13
before you say anything}. Land the printed sentence word by word: \emph{near-agreement is luck,
not evidence}. Five minutes.

\columnbreak
\textbf{We do (about 15 min).} The Guided Practice row, \emph{Our Class Survey} (problems
16--19). \textbf{This is the project, and it is live.} Three minutes to settle the item and its
four choices against the four rules (you have pre-chosen the question; the class polishes the
wording), four minutes to ask every student in the room and fill the tally, eight minutes on
16--19. \textbf{Students hold the pen; you hold the questions.}

Problem~16 goes on them cold --- population, sample, $n$, and \emph{who chose it}. Problem~17 is
the division and the class's own sentence~4. \textbf{Problem~18 is the crux transferred to the
room's own data}: if the class's top percent lands near the council's $35\%$, ask what that
proves, and make both sides defend it. Problem~19 names the limit they would fix first.

Circulate during Guided Practice and demand the reason, not the word:
\begin{itemize}[leftmargin=1.1em, nosep]
  \item \textbf{Problem 16:} ``Who is the question about? Now --- who did we actually ask?''
  \item \textbf{Problem 17:} ``Whose name is the subject of your sentence~4?'' If it is
        \emph{Riverbend students}, the design does not carry it.
  \item \textbf{Problem 18:} ``Suppose we had landed on $35\%$ exactly. Would that have made us
        right?'' A student who says yes is still at the hook vote.
  \item \textbf{Problem 19:} ``What would it cost to choose our sample by chance?''
\end{itemize}
Pick the papers to put up in the debrief --- one problem~18 that said \emph{nothing, it is luck},
and one that said the match proves the class is representative. \textbf{Your formative read comes
twice}: here, and again in the last ten minutes as students start the homework.
\end{multicols}
\end{skillbox}

\boxguard[24]
\begin{skillbox}[Debrief --- whole class, spoken (10 min)]{frost}
Pens down, papers up. \textbf{This debrief is spoken} --- there is no box at the end of the notes
to fill in, so it lives entirely on the board and in the room.
\begin{multicols}{2}
\begin{enumerate}[leftmargin=1.3em, nosep, itemsep=3pt]
  \item \textbf{Put the two panels back up} --- the council's $60$ chosen by chance, the class's
        $25$ chosen by the schedule --- and make a student say the sentence: near-agreement is
        luck, not evidence. Then state it as AFDA.DA.2j: the conclusion's \emph{validity} is
        decided by the design, not by whether the number looks plausible.
  \item Put up the two problem~18s you picked. The difference between them is not arithmetic ---
        it is whether the student still thinks matching earns the right to generalize.
  \item Put problem~9 back up and take the \emph{correct-arithmetic} reflex: \emph{``Is
        $0.36 \times 900 = 324$ wrong? Then who does $324$ describe?''} (Nobody.)
  \item Take the instrument last, from problem~5: \emph{``Don't you agree\ldots''} --- which rule,
        which 1.4 name, and why no later arithmetic repairs it.
\end{enumerate}
\columnbreak
\textbf{Check for understanding --- ask it cold, whole class.} \emph{``A homeroom of $30$ at
Bayside Middle School surveyed itself about the spirit-week theme: decades $12$, color wars $9$,
movie characters $6$, sports teams $3$. It reported to the whole school of $750$: `Bayside
students want a decades theme.' The arithmetic is right --- $12/30 = 40\%$ and
$0.40 \times 750 = 300$. Is the report honest?''} Hands down, thirty seconds of think time, then
cold-call three students. Correct answer: \textbf{no} --- the homeroom is a convenience sample, so
the honest sentence is \emph{``Of the $30$ students in this homeroom, $40\%$ chose a decades
theme''}; the $300$ describes nobody. \emph{This replaces the exit ticket.}

\textbf{Formative read.} Sort the three answers as they come: \textbf{(a)} said \emph{not honest}
and rewrote the subject; \textbf{(b)} said \emph{not honest} but could not say what to fix;
\textbf{(c)} checked the arithmetic, found it right, and called the report fine. Pile (c) is the
misconception surviving. \textbf{If pile (c) runs past a third of the room, open the unit review
by re-running \emph{The Limit} row (problems 12--15)}, and say so plainly.

\textbf{If the debrief runs short}, cut items 3 and~4 --- item~1 is the one that cannot be cut.
\textbf{Do not let the debrief eat the last ten minutes} --- the homework start is not optional
padding, it is your second formative read.
\end{multicols}
\end{skillbox}

\boxguard
\begin{skillbox}[Homework --- scored, started in class, due after two study halls]{goldbox}
Six exercises on the \textbf{Millbrook Public Library} ($1{,}200$ card holders; a phone
\textbf{SRS of $80$} against a front-desk \textbf{clipboard of $50$}): \textbf{1} the relative
frequency column ($28/80 = 35\%$, $24/80 = 30\%$, $20/80 = 25\%$, $8/80 = 10\%$); \textbf{2} read
the pre-drawn bar graph --- what a height tells a reader, and why the bars are separated (\emph{the
spiral to Lesson 1.1's variable types}); \textbf{3} the \textbf{contrast pair} --- (a)
$0.35 \times 1200 = 420$ off the SRS, (b) $30/50 = 60\%$ off the clipboard, (c) which may be
scaled and why; \textbf{4} \emph{the crux head-on} --- a clipboard showing $13$ of $50$, or
$26\%$, against the SRS's $25\%$: does the near-match earn it the school? \textbf{5} the
justification item, an instrument repair reaching back to 1.4's response bias; \textbf{6}
\textbf{the project deliverable} --- the four-sentence report on the class's own tally from the
guided notes.

\textbf{Start it in the last ten minutes of the period, in class and alone.} \textbf{Scoring:}
12 points --- 2 per item. \textbf{Item~4 is where the misconception shows}, and \textbf{item~6's
fourth sentence is what the project is graded on}: a report whose subject is \emph{Riverbend
students} has not learned the lesson, however clean its arithmetic.

\textbf{Packet or Desmos --- decided here.} The packet pages are authored and are the default.
Desmos Classroom's study-design coverage is thin, and nothing there asks a student to write a
report about their own class's data --- which is the whole deliverable --- so \textbf{1.8 is a
packet night}. Say so aloud at Close \& Assign.

\textbf{Preview.} Unit~1 ends here. Next class is the unit review, then the unit test; Unit~2
opens on one \emph{quantitative} variable, where the bar graph gives way to a dot plot and a
histogram and the bars start to touch.

\textbf{Connections \& Big Ideas.} \textit{Unit 1 core idea:} a conclusion is only as trustworthy
as the study that produced it --- run forwards, that is a design you have to defend.
\textit{Standards covered:} AFDA.DA.2d (instrument), AFDA.DA.2h (bar graph), AFDA.DA.2i (relative
frequency and scaling), AFDA.DA.2j (validity of the conclusion), PS.DC.2d (running the cycle).
\textit{Reaches back to:} 1.1's data cycle and variable types, 1.2's population and sample, 1.3's
stratified allocation, 1.4's response bias, 1.5--1.7's study types and comparisons.
\textit{Spirals forward to:} displays of quantitative data (Unit~2), and margin of error
(7.1--7.2, PS.IS.1a, d), which is the honest version of ``how close is our percent?''
\end{skillbox}

\boxguard
\begin{skillbox}[Watch For (while circulating)]{redbox}
\begin{itemize}[leftmargin=1.2em, nosep]
  \item \textbf{``We matched, so we count''} (problems 13 and~18, HW~4 --- the target
        misconception). Probe: \emph{``Who chose the people in our sample? Then what could have
        happened to our percent, and how would we know?''}
  \item \textbf{Correct arithmetic mistaken for a valid claim} (problem~9, HW~3(c)):
        $0.36 \times 900 = 324$. Probe: \emph{``Is the multiplication wrong? Then who does $324$
        describe?''}
  \item \textbf{Sentence 4 with the wrong subject} (problems 14 and~17, HW~6): ``Riverbend
        students want\ldots'' from a class of $25$. Probe: \emph{``Read me the subject of your
        sentence. Did we ask them?''}
  \item \textbf{Repairing a leading item by changing the numbers} (problems 5--6, HW~5). Probe:
        \emph{``The bias is in the wording. Which arithmetic step undoes it?''} (None.)
  \item \textbf{Overlapping choices unnoticed} (problem~6, problem~11): $14$--$16$, $16$--$18$.
        Probe: \emph{``Where does a $16$-year-old check? Now add your percents.''}
  \item \textbf{A bar's height read as a count} (problems 10, HW~2). Probe: \emph{``If the council
        had asked $120$ instead of $60$, which way does the bar move?''}
\end{itemize}
Cold-call prompts: \emph{``Name a question this class is allowed to answer about itself.''}
\quad \emph{``What would it cost us to choose $60$ students by chance?''}
\end{skillbox}

\boxguard
\begin{skillbox}[Close \& Assign (10 min)]{goldbox}
\textbf{Students start the homework here, in class, alone and silent --- this is the individual
practice, and the first minutes of it are your best formative read.} Hand the launch first ---
six exercises on paper, scored out of 12, due at the start of the first class after two study
halls --- and point out that the \emph{Keep in Mind} box on the cover is the page to flip back to,
and that \textbf{item~6 needs the tally from their own notes}, so the notes page goes home with
them. Circulate: \textbf{item~4 is the column that tells you whether \emph{The Limit} row landed},
and it is on page~2, so put students on 1, 3, and~4 first and let 2 and~5 go home; item~6 is
better written after the room has cooled. Sort what you see into three piles as you walk --- said
\emph{no} and named the design (ready); said \emph{no} with no reason (ready, needs the language);
checked the arithmetic and approved the report (the misconception survived) --- and open the unit
review by re-running problems 12--15 if that third pile ran past a third of the room. Then close
on what changed: \emph{``You walked in able to take somebody else's study apart. You walk out
having built one --- and knowing that the last sentence of your report is the part your design has
to pay for.''} \textbf{Preview:} unit review, then the Unit~1 test.
\end{skillbox}

% ── Teacher notes — one per component, in packet order (three) ──────────────
% Teacher-only prose lives HERE, never in a _key: a note is the one block in a key with no
% counterpart in the blank. No note for the debrief or the homework start — both are fully
% specified in their own boxes above.
\begin{teachernote}[Warm-Up]
Five minutes, silent, timed. Items 1 and~2 are Lesson 1.3's arithmetic on purpose --- say so, so
nobody thinks they are being retested. \textbf{Item~3 is the one to read while collecting}:
students who write $60$ and $900$ in the right rows have already made today's distinction without
the vocabulary, and rows 1 and~4 can move fast. The last question --- \emph{how many choices did
each student pick?} --- is the quiet one: \emph{exactly one} is instrument rule~3, and a student
who has it will explain problem~11 for you. Write the two counts on the board in a student's own
words and \textbf{leave them up all period}; the class's own \emph{$25$ asked, $900$ reported on}
goes under it at problem~16, and the debrief circles both lines.
\end{teachernote}

\begin{teachernote}[Guided Notes \& Practice]
\textbf{This one document is the lesson --- pace it 20 / 15, then 10 for the debrief, then 10 for
the homework start.} It is one table, and every row runs the same way: you read the printed
sentences, mark up the display, and work the first problem of the grid (1, 5, 8, 12); the class
works the rest with the pen in their hand. Rows 1 and~3 must move --- four minutes and six --- so
that row~2 gets its five and row~4 gets its five.

\textbf{Row~2 is the only new content}, and it is taught by repair: the three broken items in the
fill table are the same items Lesson 1.4 used for response bias, so the class already knows they
are wrong and is only learning what to call the rule. \textbf{Take the hand count on problem~9
before anyone writes.} If students write first, the trap is gone; if they commit to \emph{yes, the
class may report $324$} and \emph{then} are asked who chose the class, the rule arrives as the
answer to a question they already have.

\textbf{Spend the time in row~4, \emph{The Limit}}, which carries the misconception, and
\textbf{re-take the hook vote at problem~13 before you say anything}. Let the room argue. Do not
resolve it by repeating the definition --- ask who chose the class, and what would have happened
if the class had landed on $52\%$ instead.

\textbf{The Guided Practice row IS the project, and it runs on live data.} Come with the question
already chosen (the council's grant question works, and keeps the whole lesson on one context);
the class's three minutes are for polishing the wording against the four rules, not for
negotiating a topic. Then ask every student in the room --- hands, one round, four minutes --- and
fill the tally on the board while they fill theirs. \textbf{The answer key's tally is a worked
exemplar, not the right answer}: it shows a class of $25$ answering $9$ / $7$ / $6$ / $3$, which
is $36\%$ / $28\%$ / $24\%$ / $12\%$. \textbf{Your room's numbers will differ, and that is the
point} --- mark problems 16--19 and homework item~6 against your own class's totals, not against
the key's. The key's arithmetic shows the \emph{shape} of the computation; the percents are
illustrative.

If the class's top percent happens to land near the council's $35\%$, you have the lesson handed
to you --- ask problem~18 immediately. If it lands far away, that is just as good: ask whether the
council was therefore wrong. Neither outcome changes the answer.

\textbf{There is no independent practice set on this page, and that is deliberate --- the homework
is the individual practice}, and its item~6 is the project deliverable. Students need their tally
to write it, so make sure the notes page goes home.
\end{teachernote}

\begin{teachernote}[Homework]
\textbf{Scored, 12 points (2 per item), due at the start of the first class after two study
halls.} The ten minutes at the end of the period are a \emph{start}, not a completion: put
students on 1, 3, and~4 while you can still catch a wrong start, and let 2 and~5 go home.
\textbf{Item~6 is the project}, and it is the last thing to score and the first thing to read:
count sentence~4 alone --- its subject must be \emph{this class}, the sample must be named as a
convenience sample, and the limit must be stated as a fact rather than apologized for. A report
with flawless percents and \emph{``Riverbend students want\ldots''} as its conclusion earns the
arithmetic and loses the point of the lesson. \textbf{Mark item~6 against your own class's tally},
not against the key --- the key's $9$ / $7$ / $6$ / $3$ is an exemplar.

Item~4 is the fullest diagnostic in the set: a clipboard at $26\%$ beside an SRS at $25\%$, and a
student who calls that agreement \emph{evidence} is still at the hook vote. Score items 1--3 and~5
for correctness and items 4 and~6 for the \emph{reason}, not the verdict. \textbf{When you score
the page, sort item~4 into three piles} --- named the design (ready); said \emph{no} vaguely
(ready, needs the language); approved the clipboard because the arithmetic checked out (the
misconception survived). If more than a third land in the last pile, reopen the two panels in the
unit-review warm-up rather than letting it ride into the test.

\textbf{Packet is the call for 1.8}, and it is not a close one --- item~6 is a piece of writing
about data the class made this period, and no Desmos activity carries it.
\end{teachernote}
\end{document}
